% ============================================================
%  chapters/ch02-regions-boundaries-collapse.tex
% ============================================================

\chapter{Regions, Boundaries, and Collapse}

\chapterepigraph{%
  A boundary is not a wall between two things.\\
  It is the definition of both things at once.
}{Flyxion, \textit{Semantic Infrastructure}}

\sidenote{See App.\,A for\\constraint spaces;\\ App.\,C for\\Spherepop collapse}

The previous chapter established that accessibility precedes state, that
trajectories are more fundamental than things. This chapter examines the
second primitive operation from which everything else is built: \emph{containment}.

Before symbols, before operators, before logic — there is the boundary. The
inside and the outside. The contained and the containing. This observation
is not merely geometric. It is the deep structure of computation, perception,
and reasoning.

% ── 2.1 ──────────────────────────────────────────────────────
\section{Containment as a Primitive Operation}

Consider how children first encounter mathematics. Not through equations.
Through regions: \emph{this block is inside the box; that one is outside}. The
inside/outside distinction precedes counting, precedes addition, precedes
everything that adult mathematics calls fundamental. Containment is the
first cognitive operation.

\begin{definition}{Containment Relation}{containment}
  Let $X$ be a topological space. A \emph{containment relation} on
  $X$ is a binary relation $\subset$ on subsets of $X$ satisfying:
  \begin{enumerate}
    \item \textbf{Transitivity:} if $A \subset B$ and $B \subset C$
          then $A \subset C$.
    \item \textbf{Antisymmetry:} if $A \subset B$ and $B \subset A$
          then $A = B$.
    \item \textbf{Non-triviality:} there exist $A, B$ with
          $A \subsetneq B$.
  \end{enumerate}
\end{definition}

The third condition is what distinguishes a containment relation from a
trivial one. The existence of proper containment — something genuinely
inside something else — is the primitive productive operation. Everything
in logic, computation, and language follows from this.

\begin{historynote}
  George Spencer-Brown's \textit{Laws of Form} (1969) derived all of Boolean
  algebra from a single primitive: the \emph{mark}, which distinguishes an
  inside from an outside. The mark creates a boundary; the boundary creates
  two regions; all logical operations follow from the arrangement and crossing
  of boundaries. Spencer-Brown showed that the fundamental operation of logic
  is not AND or OR — it is \emph{distinction}. The Spherepop formalism can be
  understood as a dynamical extension of this insight: not just distinctions,
  but distinctions that \emph{collapse}.
\end{historynote}

% ── 2.2 ──────────────────────────────────────────────────────
\section{Nested Structures}

Containment becomes productive when it is \emph{iterated}. A region inside
a region inside a region — this is the structure of nested evaluation, nested
scope, nested context.

\begin{definition}{Nesting Sequence}{nesting-seq}
  A \emph{nesting sequence of depth $n$} is a chain
  \[
    A_1 \subsetneq A_2 \subsetneq \cdots \subsetneq A_n
  \]
  of proper containments. The \emph{depth} of an element $x \in A_1$
  within $A_n$ is $n$.
\end{definition}

Depth is the key quantity. It measures how many boundary crossings lie between
a point and the outermost context. In programming, depth is scope level.
In logic, depth is quantifier nesting. In Spherepop, depth is evaluation
priority: the innermost sphere collapses first.

\begin{example}
  In the arithmetic expression $1 + (3 \times 2^2)$, the subexpression
  $2^2$ has depth 2 (inside the multiplication, inside the addition). The
  multiplication $3 \times (\cdot)$ has depth 1. The addition has depth 0.
  Spherepop evaluation collapses innermost regions first: depth 2 before
  depth 1 before depth 0. This is not a convention — it is the \emph{only}
  consistent order for a containment-based evaluation system.
\end{example}

\begin{proposition}{Depth Determines Evaluation Order}{depth-order}
  In any containment-based evaluation system with the property that
  outer regions depend on the values of inner regions, evaluation
  must proceed from maximum depth toward minimum depth.
\end{proposition}

\begin{proof}
  Suppose an outer region $A \supset B$ is evaluated before the inner
  region $B$. Then the value of $A$ must reference $B$'s value, which
  is not yet determined — a contradiction. Therefore $B$ must be
  evaluated first.
\end{proof}

This simple observation is the entire theory of operator precedence,
scope rules, lazy evaluation, and call-by-value versus call-by-name.
All are special cases of depth-first evaluation in a containment hierarchy.

% ── 2.3 ──────────────────────────────────────────────────────
\section{Recursive Evaluation}

Depth-first evaluation in a nesting sequence is inherently recursive. The
evaluation of any region reduces to the evaluation of its interior.

\begin{definition}{Recursive Evaluation}{recursive-eval}
  Let $\mathcal{E}$ be a set of expressions with a valuation function
  $\mathrm{val} : \mathcal{E} \to V$. The valuation is \emph{recursive}
  if for any expression $e$ containing subexpressions $e_1, \ldots, e_n$:
  \[
    \mathrm{val}(e) = f(\mathrm{val}(e_1), \ldots, \mathrm{val}(e_n))
  \]
  for some combining function $f$ determined by $e$'s outermost operation.
\end{definition}

Recursion is not a programming technique. It is the mathematical structure
of evaluation in any system organized by containment. The reason recursion
feels powerful and natural — and the reason it sometimes feels dizzying —
is that it is the native language of nested regions.

\begin{conceptaside}{Recursion, Containment, and Self-Reference}
  There is a profound relationship between containment, recursion, and
  self-reference. A self-referential statement — ``this sentence is false''
  — attempts to contain itself. It is a region that includes itself as
  a subregion, violating antisymmetry of containment.

  Gödel's incompleteness theorems can be understood as a rigorous version
  of this observation: any sufficiently expressive formal system contains
  statements that refer to their own provability status. The system attempts
  to contain a description of itself within itself. The resulting
  incompleteness is not a flaw — it is the mathematical consequence of
  a sufficiently deep containment structure encountering its own boundary.

  In RSVP terms: self-reference is the collapse of a region onto its own
  accessibility set. The system attempts to be its own admissible future.
  The result is either fixed-point (stable self-reference) or oscillation
  (paradox).
\end{conceptaside}

% ── 2.4 ──────────────────────────────────────────────────────
\section{Constraint Boundaries}

Not all boundaries are walls. Some boundaries are \emph{constraint surfaces}
— the boundary of the admissible region in a constraint space.

Recall from Chapter~1 and Appendix~A that a constraint space is a triple
$(X, A, \kappa)$ where $A = \kappa^{-1}(0)$ is the admissible set. The
\emph{boundary} of this set:
\[
  \partial A = \overline{A} \setminus A^\circ
\]
is where constraint violations begin. Points inside $A^\circ$ satisfy all
constraints; points outside $A$ violate them; points on $\partial A$ are
at the exact threshold.

\begin{definition}{Constraint Boundary}{constraint-boundary}
  The \emph{constraint boundary} of a constraint space $(X, A, \kappa)$
  is $\partial A$, the topological boundary of the admissible set.
  The \emph{boundary gradient} at $p \in \partial A$ is $\nabla \kappa(p)$,
  pointing in the direction of maximum constraint violation.
\end{definition}

The constraint boundary plays the same role that a physical boundary plays
in classical mechanics: it is where dynamics change character. A trajectory
inside $A^\circ$ moves freely; near $\partial A$ it bends; crossing $\partial A$
violates the constraints. The system is reflected, refracted, or absorbed
at the boundary depending on the dynamics.

\begin{figure}[h]
  \centering
  \begin{tikzpicture}[scale=1.0]
    \tikzset{
  traj/.style={-{Stealth[length=5pt]}, thick},
  gradv/.style={-{Stealth[length=3.5pt]}, RSVPaccent!80, thin},
}
% constraint-boundary-2d.tex — constraint boundary with trajectories
  traj/.style={-{Stealth[length=5pt]}, thick},
  gradv/.style={-{Stealth[length=3.5pt]}, RSVPaccent!80, thin},
  % Admissible region (blob shape)
  \fill[RSVPlight, opacity=0.45]
    plot[smooth cycle, tension=0.7] coordinates {
      (-1.7,0) (-1.1,1.05) (0,1.45) (1.1,0.95) (1.5,0) (1.1,-0.95) (0,-1.35) (-1.05,-0.85)
    };
  \draw[RSVPdeep, very thick]
    plot[smooth cycle, tension=0.7] coordinates {
      (-1.7,0) (-1.1,1.05) (0,1.45) (1.1,0.95) (1.5,0) (1.1,-0.95) (0,-1.35) (-1.05,-0.85)
    };

  % Labels
  \node[RSVPdeep, font=\large\itshape] at (0, 0.2) {$A$};
  \node[RSVPmid, font=\footnotesize\itshape] at (0, -0.18) {$\kappa=0$};
  \node[RSVPdeep, font=\small] at (1.75, 0.55) {$\partial A$};
  \draw[RSVPdeep!50, thin] (1.58,0.52) -- (1.32,0.52);
  \node[RSVPmid!70, font=\small\itshape] at (2.1,-0.75) {$X\setminus A$};

  % Interior trajectory (free flow)
  \draw[traj, RSVPmid] (-1.1,-0.35) .. controls (-0.5,0.3) and (0.2,0.45) .. (0.75,0.1);
  \node[RSVPmid, font=\footnotesize\itshape] at (-0.2, 0.6) {$\mathbf{v}$};

  % Trajectory bending near boundary
  \draw[traj, RSVPaccent] (0.55,-0.75) .. controls (1.05,-0.45) and (1.2,0.1) .. (0.95,0.65);
  \node[RSVPaccent, font=\footnotesize\itshape] at (1.45,0.2) {$\nabla S$};

  % Collapse trajectory (outside → inside)
  \draw[RSVPred!80, thick, dashed] (1.9,-0.28) .. controls (1.75,-0.1) .. (1.35,0.2);
  \draw[traj, RSVPred!80] (1.35,0.2) .. controls (1.05,0.38) .. (0.65,0.45);
  \node[RSVPred!80, font=\footnotesize\itshape] at (2.2,-0.1) {$\mathcal{K}$};

  % Hardcoded boundary gradient arrows (pointing inward)
  \draw[gradv] (0,1.45) -- (0,1.1);
  \draw[gradv] (1.1,0.95) -- (0.85,0.72);
  \draw[gradv] (1.5,0) -- (1.15,0);
  \draw[gradv] (1.1,-0.95) -- (0.85,-0.72);
  \draw[gradv] (-1.1,1.05) -- (-0.85,0.8);
  \draw[gradv] (-1.7,0) -- (-1.3,0);

  \node[RSVPaccent!70, font=\footnotesize\itshape] at (-0.45,1.6) {$-\nabla\kappa$};

  \end{tikzpicture}
  \caption{Constraint boundaries in a two-dimensional state space.
    The admissible region $A$ (shaded) is bounded by the constraint
    surface $\partial A$ (thick curve). Trajectories inside $A^\circ$
    flow freely along the vector field $\mathbf{v}$. Near $\partial A$,
    the accessibility gradient $\nabla S$ curves trajectories inward.
    The collapse operator projects boundary-crossing trajectories
    back into $A$.}
  \label{fig:constraint-boundary}
\end{figure}

% ── 2.5 ──────────────────────────────────────────────────────
\section{Why Collapse Appears in Logic, Computation, and Perception}

The word \emph{collapse} appears throughout this monograph with a consistent
technical meaning: the reduction of a nested region to a single value,
collapsing the depth hierarchy at one level.

But collapse appears throughout mathematics and physics with the same
essential structure:

\begin{itemize}
  \item \textbf{Logical collapse:} a proof reduces a complex proposition
        to True or False. The interior of the proof — all the intermediate
        steps — collapses into a single bit.
  \item \textbf{Computational collapse:} function evaluation reduces
        a complex expression to a value. The call stack — all the nested
        frames — collapses as it unwinds.
  \item \textbf{Perceptual collapse:} visual processing reduces a
        high-dimensional retinal image to an object recognition. The
        hierarchy of features — edges, regions, parts, objects — collapses
        into a label.
  \item \textbf{Quantum collapse:} measurement reduces a superposition
        of states to a definite outcome. The probability distribution
        collapses to a point mass.
  \item \textbf{Semantic collapse:} understanding reduces a complex
        sentence to a meaning. The tree of syntactic dependencies collapses
        into an interpretation.
\end{itemize}

\sidenote{See Ch.\,8 on\\semantic\\renormalization}

These are not merely analogies. They share a common mathematical structure:
a map from a high-dimensional structured space to a lower-dimensional
value space, along the direction of decreasing depth.

\begin{definition}{Collapse Map}{collapse-map}
  Let $(X, A, \kappa)$ be a constraint space and let $\mathcal{N}(X)$
  be the set of nesting sequences in $X$. A \emph{collapse map} is a
  function
  \[
    \mathcal{K} : \mathcal{N}(X) \longrightarrow X
  \]
  satisfying $\mathcal{K}(A_1 \subsetneq \cdots \subsetneq A_n) \in A$
  (the result is always admissible) and
  $\kappa(\mathcal{K}(\cdot)) = 0$ (the result always satisfies all constraints).
\end{definition}

The collapse map is admissibility-preserving by construction. This is why
evaluation, proof, and perception all produce outputs that \emph{make sense}:
collapse maps are defined to land in the admissible region.

\begin{theorem}{Collapse Reduces Depth}{collapse-depth}
  Any collapse map $\mathcal{K}$ satisfies
  $\mathrm{depth}(\mathcal{K}(\sigma)) < \mathrm{depth}(\sigma)$
  for any non-trivial nesting sequence $\sigma$ of depth $\geq 1$.
\end{theorem}

\begin{proof}
  A non-trivial nesting sequence has depth $n \geq 1$. By definition,
  $\mathcal{K}(\sigma) \in X$ (a point, not a sequence). The depth of a
  point in $X$ is 0. Therefore $\mathrm{depth}(\mathcal{K}(\sigma)) = 0 < n$.
\end{proof}

This theorem is the reason that computation \emph{terminates}: every collapse
reduces depth by at least one, and depth is bounded below by zero.

% ── 2.6 ──────────────────────────────────────────────────────
\section{Boundaries as Information Generators}

There is a perspective on boundaries that is easy to miss: a boundary does
not divide two regions from each other — it \emph{defines} them both
simultaneously. Before the boundary, there is undifferentiated space. After
the boundary, there are two distinct regions with mutual complement structure.

\begin{proposition}{Boundary Generates Distinction}{boundary-distinction}
  Given a topological space $X$, every closed proper subset $B \subsetneq X$
  generates a distinction: the pair $(B, X \setminus B)$ of complementary
  regions. The information content of this distinction is
  $I(B) = \log_2 \frac{\mu(X)}{\mu(B) \cdot \mu(X \setminus B)}$
  for any probability measure $\mu$ on $X$.
\end{proposition}

This is why boundaries are computationally productive. A single boundary
crossing — a single pop, in Spherepop terms — generates information. The
inside and the outside become distinct. Their distinctness is what makes
further computation possible.

\begin{remark}
  The perceptron — the fundamental unit of neural networks — is a boundary
  generator. It draws a hyperplane through feature space, creating an inside
  and an outside. The entire expressive power of deep networks is built
  from iterated boundary generation. Universal approximation theorems
  (Appendix~N) say: any continuous function can be approximated by
  sufficiently many hyperplane boundaries. This is the same claim, in
  computational language, as the claim that nested containment suffices
  for all computation.
\end{remark}

% ── Exercises ────────────────────────────────────────────────
\section*{Exercises}

\begin{exercise}
  Let $X = \{1, 2, 3, 4, 5\}$ and consider the nesting sequence
  $\{3\} \subsetneq \{2,3,4\} \subsetneq X$.
  What is the depth of each element of $X$? Define a collapse map
  $\mathcal{K}$ that maps this nesting sequence to a single element
  satisfying $\kappa = 0$ where $\kappa(x) = |x - 3|$.
\end{exercise}

\begin{exercise}
  Formalize the claim that depth-first evaluation is the only consistent
  order for containment-based evaluation. Specifically: suppose evaluation
  proceeds by some other order (breadth-first, random). Construct an
  expression for which this order produces an undefined or contradictory
  result.
\end{exercise}

\begin{exercise}[title={Topological}]
  The boundary $\partial A$ of a closed set $A$ satisfies
  $\partial A = \overline{A} \cap \overline{X \setminus A}$. Prove that
  $\partial A$ is always closed. Give an example where $A$ is open,
  $A$ is closed, and $A$ is neither open nor closed, showing what
  $\partial A$ looks like in each case.
\end{exercise}

\begin{exercise}[title={Open problem}]
  Define an appropriate notion of \emph{collapse entropy}: the information
  lost when a nesting sequence of depth $n$ is collapsed to a point.
  Under what conditions is collapse entropy minimal? Does minimal-entropy
  collapse correspond to any known computational strategy (memoization,
  normal-form reduction, canonical representatives)?
\end{exercise}
